\chapter{从数据到推理的完整流程}
\label{chap:pipeline}
\begin{quote}\small
\textit{本章摘要。}本章把前面定义的学习问题与经典模型落实为可执行流程：先定义事件和样本，再划分数据、训练模型、进行推理并评价部署风险。数据治理和时间边界贯穿全章，因为任何后续算法都无法修复样本构造或信息泄漏造成的错误。
\end{quote}
前一章说明了模型如何依据目标求解，却有意把设计矩阵当作已经给定。工业任务最容易出错的地方，恰恰发生在矩阵形成之前：一条记录何时进入系统，一个窗口对应哪个未来事件，哪些设备可以出现在训练和测试两侧。这些选择不是附属工程细节，而是在决定风险公式中的样本究竟来自什么分布。

本章会反复检查一个朴素的问题：作出预测的那一刻，这条信息真的已经到手了吗？标准化的均值从哪里算，参数用哪些记录训练，阈值在哪一批数据上选择，都要留下清楚的记录。评价时还要把相邻窗口合并回故障事件：同一次故障连续触发许多条告警，并不等于发现了许多次故障。否则离线分数可能很好看，现场却仍然漏报或频繁打扰检修人员。

\section{数据划分与数据治理}
数据泄漏不仅发生在训练测试样本重叠时，也可能通过预处理、特征选择与标签构造进入整个流程\cite{kaufman2012leakage}。因此以下划分规则约束的是信息使用过程，而不只是一次随机切分。
训练集用于拟合参数，验证集用于选择超参数和比较方案，测试集只在最终报告时使用。若时间顺序具有因果意义，应采用按时间切分，而不是随机打乱。工业数据还需处理缺失、漂移、传感器失真、重复记录与标签延迟。数据清洗不是模型之外的琐事，因为它决定了模型究竟看到了什么分布。

\paragraph{预处理也是需要拟合的模型。}
设训练索引为$\mathcal T$，第$j$个特征的均值和尺度应只在训练集计算：
\begin{equation}
 \hat\mu_j=\frac1{|\mathcal T|}\sum_{i\in\mathcal T}x_{ij},\qquad
 \hat s_j^2=\frac1{|\mathcal T|}\sum_{i\in\mathcal T}(x_{ij}-\hat\mu_j)^2,
 \qquad z_{ij}=\frac{x_{ij}-\hat\mu_j}{\max(\hat s_j,\epsilon)}.
\end{equation}
这里$\epsilon>0$是与该特征单位一致的尺度下限，避免常数列除零；也可明确删除常数列。验证、测试与上线输入都使用同一组训练统计量。其原因可写为函数依赖：最终预测是$f_{\hat\theta}(T_{\hat\phi}(x))$，其中$\hat\phi=\operatorname{FitTransform}(\mathcal D_{\mathcal T})$。若用全数据估计$\hat\phi$，测试集就参与了模型构造，即使没有使用测试标签，也改变了原定的归纳评估问题。

交叉验证时每折必须重新拟合$\hat\phi$与$\hat\theta$，而不是先全量标准化再切折。缺失值填补、PCA、异常值阈值和监督特征筛选也属于这条规则；若部署允许使用一批无标签目标数据适配，应将其明确写成另一个评估规则。

\section{训练、评价、部署与反馈}
\begin{figure}[htbp]
\centering
\begin{tikzpicture}[node distance=0.7cm and 0.65cm]
 \node[idi input, minimum width=2.0cm] (raw) {原始记录\\时间戳、设备};
 \node[idi process, right=of raw, minimum width=2.1cm] (window) {窗口与标签\\只用可用信息};
 \node[idi state, right=of window, minimum width=2.0cm] (split) {按时间/设备\\划分数据};
 \node[idi process, right=of split, minimum width=2.0cm] (fit) {训练与验证\\模型版本};
 \node[idi output, right=of fit, minimum width=2.0cm] (serve) {上线推理\\告警与回退};
 \draw[idi arrow] (raw)--(window); \draw[idi arrow] (window)--(split);
 \draw[idi arrow] (split)--(fit); \draw[idi arrow] (fit)--(serve);
 \draw[idi feedback] (serve.north) .. controls +(0,0.68) and +(0,0.68) .. node[above,font=\scriptsize]{延迟标签与审计} (window.north);
 \draw[idi feedback] (split.south) .. controls +(0,-0.42) and +(0,-0.42) .. node[below,font=\scriptsize]{测试集只用于最终报告} (fit.south);
\end{tikzpicture}
\caption{工业学习管道中的信息边界。切分、特征生成和标签构造共同决定是否发生泄漏。}
\label{fig:pipeline-information-boundary}
\end{figure}
图\ref{fig:pipeline-information-boundary}把窗口构造、时间与设备划分和上线反馈连接起来；其中延迟标签返回样本构造环节，而测试集边界限制了模型选择可使用的信息。
\paragraph{贯穿案例的固定设置。}
先以“预测未来二十四小时内的设备故障”建立维护流程的预测任务。设备$i$在预测时刻$t$的可用观测记为
\begin{equation}
 \mathcal{I}_{i,t}=\{x_{i,s}:s\leq t\}\cup\{m_{i,s}:s\leq t\},
\end{equation}
其中$s$表示记录进入系统并可查询的时刻，原始采集时刻另行保存；$x_{i,s}$表示传感器和控制量，$m_{i,s}$表示维护信息。二十四小时故障标签定义为
\begin{equation}
 y_{i,t}=\mathbb{1}\{\text{设备 }i\text{ 在 }(t,t+24\mathrm{h}]\text{ 内发生经确认的故障事件}\}.
\end{equation}
这里的预测时刻、事件发生时刻和记录可用时刻必须分开保存。计划保养并不自动构成故障标签，维修工单只是确认事件的证据来源。若未来二十四小时尚未完整观测，且还没有确认故障，该样本不能直接标为正常。二十四小时是预测时间窗，实际提前量则是故障发生时刻与有效告警时刻之差，需要单独评价。后文的剩余寿命回归、图像定位和维修问答是关联任务，各有标签与指标；共享的是证据可用性原则，而非强行统一输出。

后文比较模型时，都沿用这项任务定义，并逐项记录输入、结构、损失和测试划分。还要写清分数超过多少才告警，传感器缺失时是沿用备用规则还是请人工检查。例如，模型能算出故障概率，并不意味着它已经规定了谁来接收告警、何时安排复检。把这些安排与模型结果分别记录，才能复查成绩差异来自算法还是处理流程。

\paragraph{可用性、标签成熟与时间隔离的形式化。}
对记录$r$分别保存采集时间$e(r)$和可用时间$a(r)$。预测时刻$t$、回看长度$w>0$对应的输入集合应为
\begin{equation}
 \mathcal R_t=\{r:t-w\leq e(r)\leq t,\ a(r)\leq t\},\qquad
 x_t=\operatorname{Aggregate}(\mathcal R_t).
\end{equation}
因此仅比较采集时间不够；迟到数据不能反向修正历史在线输入而仍声称是实时评估。可将这一原则理解为$x_t$必须由$t$时刻已经掌握的信息计算出来。

令$h=24\mathrm h$，训练快照截止于$T$。若标签确认最多延迟$\delta$，且事件追踪完整，则要求$t+h+\delta\leq T$可保守保证一个二元标签成熟。已确认的早期正例虽可提前得知，但不能把尚未成熟的其余样本直接当负例；只保留这些早期正例又会改变样本分布。若延迟没有确定上界，应该依据逐样本的标签成熟状态筛选，或使用明确处理删失的模型。

若训练预测时刻为$t_a$，验证最早预测时刻为$t_b$，为使训练标签不跨入验证预测期，可要求$t_a+h<t_b$。若还要求原始窗口与标签支撑区间完全不重叠，则由$[t_a-w,t_a+h]$与$[t_b-w,t_b+h]$分离得到更保守条件$t_a+h<t_b-w$。后一条件不是所有滚动预测都必须采用：部署时合理地重用过去观测并不自动泄漏，但会造成统计依赖；应根据要模拟的部署方式选择规则并说明残留相关性。

\paragraph{优化：如何找到参数。}
以可微目标$J(\bm{\theta})$为例，梯度下降更新为
\begin{equation}
 \bm{\theta}_{t+1}=\bm{\theta}_t-\eta_t\nabla_{\bm{\theta}}J(\bm{\theta}_t),
\end{equation}
其中$\eta_t$是学习率。随机梯度下降用小批量$B_t$近似全量梯度：
\begin{equation}
 \bm{\theta}_{t+1}=\bm{\theta}_t-\eta_t\frac{1}{|B_t|}\sum_{i\in B_t}\nabla_{\bm{\theta}}\ell_i(\bm{\theta}_t).
\end{equation}
动量、Adaptive Moment Estimation (Adam)和学习率调度都改变了优化轨迹，但不改变“依据目标函数更新参数”的基本逻辑。优化收敛并不自动意味着泛化良好，因为训练目标只是有限数据上的代理目标。

\paragraph{推理与部署。}
训练完成后冻结参数，在新样本上计算$f_{\hat{\bm{\theta}}}(\bm{x})$或$p_{\hat{\bm{\theta}}}(y\mid\bm{x})$。部署阶段还要考虑延迟、吞吐、内存、量化误差、异常输入和模型更新策略。一个在离线测试集上准确的模型，如果无法在传感器数据到达后的时间约束内完成推理，仍然不是有效的工业系统。

\paragraph{评价指标。}
分类可报告准确率、精确率、召回率、$F_1$和Area Under the Receiver Operating Characteristic Curve (ROC-AUC)；回归常用Mean Absolute Error (MAE)与Root Mean Squared Error (RMSE)。概率预测还应检查负对数似然和校准误差。指标必须对应决策目标。例如，安全监测通常优先控制漏报，故不能只使用准确率。报告结果时应同时给出数据划分、随机种子、置信区间或多次运行波动，以区分真实改进和偶然差异。

\paragraph{从混淆矩阵推导指标及其基率依赖。}
固定阈值后，真正例、假正例、假负例与真负例数量记为$TP,FP,FN,TN$。于是
\begin{equation}
 \operatorname{Precision}=\frac{TP}{TP+FP},\quad
 \operatorname{Recall}=\frac{TP}{TP+FN},\quad
 F_1=\frac{2TP}{2TP+FP+FN}.\label{eq:pipeline-classification-metrics}
\end{equation}
式\eqref{eq:pipeline-classification-metrics}中的$F_1$由精确率与召回率的调和平均得到；分母为零时必须预先约定报告方式，不能把无正例设备的指标静默设成满分。

令部署故障比例为$\pi=P(Y=1)$，真正例率为$u=P(\hat Y=1\mid Y=1)$，假正例率为$v=P(\hat Y=1\mid Y=0)$。由条件概率和全概率公式，
\begin{equation}
 P(Y=1\mid\hat Y=1)=\frac{\pi u}{\pi u+(1-\pi)v}.
\end{equation}
例如$\pi=0.01,u=0.9,v=0.05$时，精确率约为$0.154$：高召回仍可能伴随大量无效告警。因此在故障比例不同的设备群间，不能只依据同一受试者工作特征（Receiver Operating Characteristic，ROC）工作点推断告警质量。ROC-AUC考察排序，部署动作还需要固定阈值、基率与成本。

\paragraph{窗口误差如何转换为事件指标。}
设第$e$次故障时间为$\tau_e$，要求最小提前量为$L_{\min}$，其中$0<L_{\min}\leq h$。只有位于$[\tau_e-h,\tau_e-L_{\min}]$的有效告警才算及时命中。定义$a_e$为该区间内按既定合并规则选取的首个有效告警，事件召回率为
\begin{equation}
 \frac{\sum_e\mathbb1\{a_e\text{存在}\}}{\text{可评价故障事件数}},\qquad
 L_e=\tau_e-a_e.
\end{equation}
未命中事件不能从召回分母移除；仅在命中事件上报告平均提前量也应明确这一条件。必须事先规定同一告警能否匹配多个事件、连续超阈值如何合并、告警冷却时间如何设置。无效告警频率宜用“未匹配告警数/有效监测设备天数”计算，而不是把重复窗口当成独立动作。

\paragraph{误差条为何依赖抽样单位。}
对固定模型的$n$个独立测试错误指示$E_i\in\{0,1\}$，若错误率为$q$，则$\hat q=n^{-1}\sum_iE_i$且$\Var(\hat q)=q(1-q)/n$，标准误可用$\sqrt{\hat q(1-\hat q)/n}$估计；小样本或极端比例应使用更合适的区间方法。若同一设备有$m$个窗口，设备间独立、设备内错误指示方差为$\sigma^2$且两两相关系数为$\rho$，共$G$台设备，则
\begin{equation}
 \Var(\bar E)=\frac{Gm\sigma^2+Gm(m-1)\rho\sigma^2}{(Gm)^2}
 =\frac{\sigma^2}{Gm}\{1+(m-1)\rho\}.\label{eq:pipeline-cluster-variance}
\end{equation}
式\eqref{eq:pipeline-cluster-variance}说明窗口数巨大不代表独立证据同样巨大。可按设备进行成组重采样，时间相关则使用适当的块重采样；重采样方案仍需匹配目标设备总体及依赖结构。多随机种子波动主要衡量训练随机性，不能替代测试抽样的不确定性。

\begin{remark}
 一条实用的诊断顺序是：先检查标签和切分，再检查训练与验证曲线，随后检查错误样本，最后才调整模型结构。否则，复杂模型很容易被用来掩盖数据问题。
\end{remark}

\paragraph{从原始记录到学习样本。}
工业数据通常不是天然的$(x,y)$对，而是多个系统中异步产生的事件。构造样本需要确定实体键、观测时间、预测时间、标签时间和允许使用的信息集合。对预测时刻$t$，输入可写成
\begin{equation}
 x_t=\operatorname{Aggregate}\{r_s:s\in[t-w,t]\},
\end{equation}
标签则由部署时真正关心的未来区间生成。若把$t$之后才产生的维修结果加入$x_t$，离线指标会被人为抬高。

落实到日志时，输入必须同时满足“已经采集”和“已经可用”。维修工单可能在故障发生数小时后才录入，可以事后用于确认标签，却不能被反向放入预测时刻的输入特征。标签审计还要区分故障维修和计划保养，避免模型实际学习维护排班而非故障风险。

该问题至少需要三组切分：按时间向后滚动的验证集、完全留出的后期测试集，以及按设备留出的迁移测试集。训练阶段可以使用窗口统计、频谱特征和最近一次维护信息，但每个特征都要通过时间可用性检查。常见泄漏包括用全周期均值标准化、用故障后的完整轨迹计算斜率、把同一设备的相邻窗口随机分到不同集合，以及使用在全数据上拟合的异常阈值。每项检查都应有一个可失败的单元测试或样本审计报告。

上线验收不应只看 ROC-AUC。系统还要规定每台设备每天最多允许的告警数、最小提前量、连续超阈值窗口数和人工确认时限。传感器断流时，服务应返回“数据不足”状态并触发数据质量工单；模型超时则回退到上一版经过验证的模型或规则基线。这样，模型风险、数据风险和服务风险才能在运行记录中被区分，而不是都表现为一个模糊的“预测错误”。

\paragraph{训练、验证和测试的职责。}
训练集回答“参数如何拟合”；验证集回答“在候选方案中如何选择”；测试集回答“在选择完成后能否泛化”。当研究者根据测试集反复改模型时，测试集已经承担验证集职责，最终数字会有乐观偏差。时间序列还应使用滚动验证，设备任务则应按设备分组，避免同一实体的相邻窗口跨越划分边界。

\paragraph{数据管道的可追溯性。}
每个特征都应能追溯到原始字段、单位、时间窗口、清洗规则和版本。保存最终矩阵还不够，因为矩阵无法说明它是否使用了未来统计量。建议同时保存样本 Identifier (ID)、设备 ID、预测时间、标签版本、特征生成版本和划分哈希。这样错误样本才可以从模型输出反查到数据源。

\paragraph{评价不等于排序。}
准确率和 RMSE 适合比较某些预测任务，却不直接等于业务价值。一个故障检测器可能提高宏平均$F_1$，却增加大量无效告警；一个回归器可能降低平均误差，却在关键峰值时系统性低估。指标应由决策流程倒推，并报告阈值、事件合并规则、提前量、资源预算和错误成本。

\paragraph{上线前后的差异。}
离线评估使用固定数据和批量计算，线上系统还会遇到数据延迟、网络断开、字段缺失、模型版本不一致和标签延迟。部署验收应至少包括正常输入、缺失输入、极值输入、重复输入、乱序输入和服务超时。模型服务失败时必须有明确回退策略，而不是返回一个未校验的默认预测。

\paragraph{本章小结。}
完整流程不是“下载数据、训练模型、报告分数”，而是把现实过程转换为可复现的学习问题，再把模型输出转换为可审计的决策。数据定义、划分规则、损失函数、优化过程、评价指标和部署约束必须互相一致；其中任何一环含糊，后续的精度数字都难以解释。

本章最重要的边界有两条：一是信息边界，任何输入只能使用预测当时已经可获得的证据；二是评价边界，训练、选择和最终检验不能共用同一份证据而不承认选择偏差。标准化统计量、标签成熟规则、事件合并与置信区间都应服从这些边界，不能仅把模型权重当作完整实验。

确定样本、标签和数据用途之后，下一章开始逐步更新参数：先沿计算图求梯度，再检查步长、曲率、初始化、归一化与数值精度。优化器可以帮助降低已经定义好的损失，却不会发现工单录入晚了，也不会替我们纠正训练集和测试集的混用。